\documentclass{article}
\usepackage{amssymb}
\usepackage{amsmath}
\usepackage[margin=1in]{geometry}
\usepackage{enumitem}
\usepackage{color}

\newcommand{\rd}{\mathrm{d}}

\title{Peer Review\\[2pt]
\large\emph{Derivation of the Analytic Solution of the Coupled DL--AL Problem
via the Time--Laplace / Resolvent Construction}}
\author{Referee report}
\date{June 2026}

\begin{document}
\maketitle

\section*{Recommendation}
\textbf{Major revision.} The overall strategy (Laplace transform $\to$ resolvent
$\to$ pole spectrum) is correct and well motivated, the algebra I checked is
sound, and the validation section is commendably honest. However, there is a
substantive disconnect between the construction \emph{derived} (Secs.~5--8) and
the one actually \emph{validated} (Sec.~9), together with an internal
inconsistency in the pole/residue bookkeeping that bears directly on the central
``analytic solution'' claim. Both are fixable, largely by reframing rather than
re-deriving.

\section{Summary of the manuscript}
The paper studies a coupled diffusive--layer (DL) / advective--layer (AL) system
on the unit disk with $N$ evenly spaced microtubules. It (i) Laplace--transforms
in time and eliminates the AL density exactly as a linear functional of the DL
trace on a microtubule ray; (ii) rewrites the microtubule sources as a Dirac comb
so they feed every angular harmonic; (iii) constructs a closed modified--Bessel
radial Green's function; (iv) closes the problem into a Fredholm equation of the
second kind whose secular determinant $\det(\mathcal I + \mathcal A(s))=0$ yields
the (complex) decay rates; (v) inverts by the Bromwich contour to a modal
residue series; and (vi) gives a two--mode truncation and a numerical validation
against a finite--difference solver across the $(v,w)$ parameter plane.

\section{Strengths}
\begin{itemize}[leftmargin=1.4em]
\item The motivation for a resolvent (rather than eigenfunction) construction on a
non--self--adjoint operator (Sec.~1) is correct and clearly argued.
\item The exact AL elimination [Eq.~(6)] is clean and correct.
\item The Dirac--comb, full--disk reformulation (Sec.~3) is elegant and is the
right way to dissolve the apparent ``extra boundary condition'' / cusp paradox of
the single--wedge treatment.
\item The radial Green's function [Eq.~(13)] is correct; I verified the Abel
constant and the central--$\delta$ collapse to the closed form [Eq.~(15)].
\item The modal--amplitude formula [Eq.~(22)] via a biorthogonal projection is
correct, and its reduction to the classical Fourier--Bessel coefficient in the
uncoupled limit (Sec.~7 sanity check) is convincing.
\item The validation (Sec.~9) is transparent and, unusually, reports a regime in
which the method fails rather than only its successes.
\end{itemize}

\section{Major concerns}

\paragraph{M1. The derived object (Secs.~5--8) is not the object validated
(Sec.~9).}
Sections 5--8 present the analytic solution as the resolvent / secular determinant
$\det(\mathcal I + \mathcal A(s)) = 0$, which is a \emph{nonlinear} eigenvalue
problem: $\mathcal A$ depends on $s$ through both the Green's kernel and the AL
functional. Section 9, however, states that the modes are obtained by writing the
coupled operator in the Bessel basis and \emph{diagonalizing} it --- an ordinary,
$s$--independent (generalized) eigenproblem for the generator. These are related
but not identical constructions, and the manuscript silently substitutes one for
the other while continuing to cite Eqs.~(24)--(25) as though they were what is
computed. The paper should either (a) prove the equivalence (generator
eigenvalues $=$ roots of $\det(\mathcal I+\mathcal A(s))=0$), or, preferably,
(b) promote the generator eigenproblem to the primary construction --- it is what
actually runs and is validated --- and present the resolvent/secular determinant
as motivation. As written, Sec.~9 cannot be reproduced from Secs.~5--8.

\paragraph{M2. The ``subdominant $\Phi_0$'' argument (Sec.~6) is internally
inconsistent and contradicts Sec.~9.}
In assembling the residue series, the manuscript discards the forced--response
residue family at the bare poles $s=-\kappa_{0,j}^2$, describing it as a
``diffusive--background residue family \dots\ which we suppress here as
subdominant.'' But the slowest bare pole, $-\kappa_{0,1}^2\approx-5.78$, is the
\emph{most} dominant (slowest--decaying) contribution, not a subdominant one.
Moreover, the dominant decay rate reported throughout Sec.~9 is
$\sigma\approx 5.0\text{--}5.8$ --- i.e.\ precisely the (perturbed) $n=0$ pole that
Sec.~6 suppresses. The text must resolve whether the bare pole is (a) removable in
$(\mathcal I+\mathcal A)^{-1}\Phi_0$, (b) shifted by the coupling (as claimed near
Eq.~(13)), or (c) a genuinely separate contribution; at present all three are
asserted at different points. This is the single most important rigor gap, since
it concerns which poles constitute ``the spectrum.''

\paragraph{M3. The ``truncation breakdown'' diagnosis is asserted, not
demonstrated.}
For the failing regime $v=100,\ w=100$ (predicted $\sigma=11.34$ vs.\ numerical
$3.40$; mass MSE $\sim10^{-1}$), the paper attributes the disagreement to a
too--coarse few--mode truncation together with first--order discretization
effects. No convergence evidence is given: it is never shown that increasing the
number of radial/angular modes (or refining the grid) reduces the error
monotonically. Without such a study the failure could equally reflect a genuine
model mismatch. A mode--count and grid--refinement convergence panel for this
regime should be added before the ``truncation'' explanation can be accepted.

\paragraph{M4. The validated equations are not the continuum equations of
Secs.~1--2.}
Section 9 concedes that the analytic operator is supplied with the discrete
scheme's coupling, $a \to a\,r\,\Delta\theta$ --- a grid--dependent weight with
\emph{no continuum limit}. Consequently Sec.~9 validates the analytic solution of
a grid--modified system, not the continuum problem derived in Secs.~1--8. This
caveat should be stated at the \emph{opening} of Sec.~9 (and reflected in the
abstract), not introduced midway. Relatedly, the phrase ``depend \emph{mildly} on
the grid'' understates the effect: the paper's own refinement data move $\sigma$
by roughly $25\%$ ($\approx4.2\to5.2$) over $16\to48$ rings. ``Mildly'' should be
removed or quantified.

\paragraph{M5. Non--self--adjoint completeness is assumed silently.}
The modal series [Eq.~(24)] presumes simple poles, the absence of Jordan blocks,
and completeness of the biorthogonal family $\{\psi_k\}$. For a non--self--adjoint
operator none of these is automatic. At minimum the manuscript should add an
explicit caveat; ideally it should argue (or cite) why the expansion is complete
for this operator.

\section{Minor concerns}
\begin{enumerate}[leftmargin=1.6em,label=m\arabic*.]
\item \textbf{Notation clash.} $a$ denotes the on--rate throughout, yet Eq.~(27)
writes the angular--modulation amplitude as $a(r)$. Rename the amplitude (e.g.\
$\alpha(r)$) to avoid collision with the coupling rate.
\item \textbf{Addition--theorem hand--wave (Eq.~(18)).} The remark that the
ray--summed kernel ``is the image--corrected $\tfrac{1}{2\pi}K_0(\sqrt s|\mathbf
r-\mathbf r'|)$'' is loose: the sum runs over orders $nN$ (not all $n$), and the
disk Dirichlet correction is not the free--space $K_0$. Either derive the image
series or delete the parenthetical and retain $\mathcal S$ as the mode sum.
\item \textbf{Undefined block structure.} In the $2\times2$ truncation [Eq.~(30)],
$\mathcal S_{mn}$ is called ``the $(m,n)$ angular block of $\mathcal S$,'' but
$\mathcal S$ [Eq.~(18)] is defined as a single angle--summed scalar kernel with no
$(m,n)$ block structure. The projection producing $p_{mn},q_{mn}$ is asserted
rather than derived.
\item \textbf{Sign of $v$.} The continuum AL equation uses $+v\,\partial_r\rho$;
the physical (inward) direction and the sign convention for $v$ are never fixed.
State them explicitly (the companion solver negates the velocity input).
\item \textbf{Branch cut.} The dismissal of the $\sqrt s$ cut (Sec.~6) is correct
for a bounded domain but would be clearer with a one--line note that
$\Phi_0,\mathcal K_m$ are even in $\sqrt s$ and hence meromorphic in $s$.
\item \textbf{Title.} Given that the secular roots are only obtained numerically
and the validated object is grid--modified, ``Derivation of the Analytic
Solution'' slightly oversells; ``A resolvent construction\dots'' or
``semi--analytic'' would be more accurate.
\item \textbf{Reproducibility.} The exact \texttt{deterministic\_verification}
invocation (parameters, mode counts, time stamps) that produced Table~1 should be
recorded for the notes' own future use.
\end{enumerate}

\section{Questions to the authors}
\begin{enumerate}[leftmargin=1.6em]
\item Are the generator eigenvalues of Sec.~9 provably the roots of
$\det(\mathcal I+\mathcal A(s))=0$ of Sec.~5? If so, the equivalence should be
shown; if not, which set constitutes the physical spectrum?
\item Does refining the mode counts $n_{\max},J_r$ and the grid reduce the
$v=100,\ w=100$ error monotonically?
\item Since the dominant root $s_k$ lies near a bare pole $-\kappa_{0,1}^2$, is the
kernel $\mathcal K_0(\,\cdot\,;s_k)$ not nearly singular there? How is this handled
numerically in the reconstruction [Eq.~(24)]?
\end{enumerate}

\section{Assessment}
The physical model and the bulk of the algebra are correct, the comb
reformulation is a genuine contribution, and the honest validation (including a
failing regime) is a real asset. The required revisions are primarily about
\emph{coherence}: align the derived construction (Secs.~5--8) with the validated
one (Sec.~9) [M1], repair the pole/residue bookkeeping [M2], and foreground the
fact that the validated system is the grid--modified one [M4]. The one outstanding
\emph{empirical} task is a convergence study for the failing regime [M3]. With
these addressed, the manuscript would constitute a sound and useful account of the
coupled DL--AL problem.

\end{document}
